\documentclass[10pt]{article}
\usepackage[utf8]{inputenc}
\usepackage[margin=0.75in]{geometry}
\usepackage{amsmath,amssymb}
\usepackage{hyperref}
\usepackage{microtype}

\hypersetup{
    colorlinks=true,
    linkcolor=blue,
    citecolor=blue,
    urlcolor=blue
}

\begin{document}
\pagestyle{empty}

\noindent
\textbf{Felipe Mora, M.Sc.} \\
Departamento de Industrias, Universidad Técnica Federico Santa María \\
Av. España 1680, Valparaíso, Chile \\
Email: \href{mailto:felipe.morar@usm.cl}{felipe.morar@usm.cl} \quad | \quad ORCID: \href{https://orcid.org/0009-0001-1034-5948}{0009-0001-1034-5948} \\[2.5mm]

\noindent
\today \\[2mm]

\noindent
\textbf{The Editors-in-Chief} \\
\emph{International Review of Financial Analysis} \quad (Elsevier) \\[2mm]

\noindent
\textbf{Re: Submission of Original Research Article ``Causal Emergence in U.S. Equity Industry Portfolios: Dynamic Organization and Effective Dimensionality During Systemic Stress''} \\[2mm]

\noindent
Dear Editors-in-Chief,

\vspace{1mm}
\noindent
I am pleased to submit my original research manuscript titled \textbf{``Causal Emergence in U.S. Equity Industry Portfolios: Dynamic Organization and Effective Dimensionality During Systemic Stress''} for publication consideration in the \emph{International Review of Financial Analysis} (IRFA).

\vspace{1mm}
\noindent
A central question in empirical financial economics and systemic risk analysis is whether collective financial market distress develops an organized macroscopic dynamical structure that cannot be reduced to contemporaneous covariance, common factor exposure, or univariate persistence.

\vspace{1mm}
\noindent
In this study, I operationalize continuous-state \emph{Causal Emergence} (CE) across daily U.S. equity industry portfolio returns from January 1990 to June 2026. By modeling cross-sector transitions via Vector Autoregressions and optimizing coarse-graining projection matrices on Stiefel manifolds under scale-adaptive Gaussian interventions, I introduce: (1) the \textbf{Causal Emergence Financial Index} ($\mathrm{CEFI}_t$), quantifying macroscopic effective information density gained by coarse-graining (in nats/DOF); and (2) the \textbf{Causal Effective Dimension} ($q_t^*$), identifying the optimizer-selected macro dimension concentrating dynamic transitions.

\vspace{1mm}
\noindent
To account for finite-sample optimization effects, I evaluate historical estimates against a four-tier hierarchy of matched surrogate null models ($B=9,999$ for primary nulls). Key empirical findings include:
\begin{itemize}\setlength{\itemsep}{0pt}\setlength{\parskip}{0pt}\setlength{\parsep}{0pt}
    \item \textbf{Descriptive Regime Variation:} $\mathrm{CEFI}_t$ varies across historical stress episodes. The descriptive excess over cross-lag network isolation is larger at the 2008 GFC peak (+0.0778 vs. +0.0370 nats/DOF), while the 2005 calm benchmark is statistically consistent with surrogate persistence ($p_{\text{Holm}} \ge 0.3480$).
    \item \textbf{Matched Surrogate Inference:} After Holm-Bonferroni correction, no primary comparison rejects at the 5\% FWER ($p_{\text{Holm}} \ge 0.1092$). While the March 2020 COVID shock nominally exceeds static correlation modes ($p_{\text{emp}} = 0.0182$), it remains consistent with cross-lag isolation ($p_{\text{emp}} = 0.2893$).
    \item \textbf{Conditional Low Dimensionality:} Under baseline Gaussian interventions ($\kappa = 1.0$), 98.94\% of rolling windows select $q^* \le 4$ (modal $q^* = 1$ accounting for 63.8\%); matched surrogate nulls also select modal $q_0^* = 1$, indicating that low dimensionality is also favored by the per-degree-of-freedom criterion under matched surrogate models, and thus is not, by itself, evidence of non-null emergent organization.
\end{itemize}

\vspace{1mm}
\noindent
The manuscript has not been published previously, is not under consideration elsewhere, and its submission is approved by the author. I confirm compliance with all ethical and editorial policies of the journal.

\vspace{2.5mm}
\noindent
Sincerely, \\[1.5mm]
\textbf{Felipe Mora, M.Sc.} \\
Departamento de Industrias, Universidad Técnica Federico Santa María \\
Email: \href{mailto:felipe.morar@usm.cl}{felipe.morar@usm.cl}

\end{document}
